\documentclass[11pt]{article}

\usepackage[margin=1in]{geometry}
\usepackage{amsmath,amssymb,amsthm}
\usepackage{graphicx}
\usepackage{hyperref}
\usepackage{xcolor}
\hypersetup{hidelinks}

\newtheorem{theorem}{Theorem}
\newtheorem{lemma}[theorem]{Lemma}
\newtheorem{corollary}[theorem]{Corollary}
\newcommand{\F}{\mathbb F}
\newcommand{\E}{\mathbb E}
\newcommand{\one}{\mathbf 1}
\newcommand{\Nphi}{\mathcal N_{\Phi}}

\title{Optimality of the endpoint Orlicz loss:\\
every dense Walsh subsystem contains an affine spike}
\author{Run \texttt{fa\_banach\_001}, lane 19}
\date{August 11, 2026}

\begin{document}
\maketitle

\begin{abstract}
Burstein proves that every uniformly bounded orthogonal system of size $n$
has a subsystem of cardinality at least
$n/[e(\log n)^{\alpha+1}]$ whose $L^2\log^\alpha L$ synthesis norm is at
most $C_\alpha(\log\log n)^{\alpha/2}$, and asks whether the iterated
logarithm is optimal.  We prove that it is.  For every sufficiently large
$N=2^d$, take all characters of $\F_2^d$.  Every subset of the target
cardinality contains an affine frequency subspace of size
$m\asymp_\alpha \log N/\log\log N$.  The normalized character sum on that
subspace is exactly $\sqrt m$ on a set of measure $1/m$ and zero elsewhere,
so its Luxemburg norm is at least
$c_\alpha(\log m)^{\alpha/2}\asymp
c_\alpha(\log\log N)^{\alpha/2}$.  The affine-subspace assertion follows
from an elementary cube count and a count of linearly dependent direction
tuples.  Thus the source upper bound has the optimal asymptotic order in the
worst case.
\end{abstract}

\section{Source question}

The source is Will Burstein, \emph{$\Lambda_p$ Style Bounds in Orlicz Spaces
Close to $L^2$}, arXiv:2505.00155v2 \cite{Burstein}.  Its Theorem 2.7 proves
the upper bound discussed below.  On page 13, Remark 2.11 says that the
optimality of equation (2.4) remains open.  The relevant source passage is
reproduced here.

\begin{center}
\includegraphics[width=0.98\textwidth]{figures/open_problem_crop.png}
\end{center}

Fix $\alpha>0$.  The source uses a Young function $\Phi$ for which, with
constants $u_0\ge e$ and $c(\alpha)>0$ fixed once $\alpha$ is fixed,
\[
 \Phi(u)=u^2\log^\alpha u\quad (u\ge u_0),
 \qquad
 \Phi(u)=c(\alpha)u^2\quad (0\le u<u_0).
\]
The corresponding Luxemburg norm on a probability space is
\[
 \|f\|_\Phi
 =\inf\left\{\lambda>0:
 \E\Phi(|f|/\lambda)\le1\right\}.
\]
For an orthogonal family $(\varphi_\xi)_{\xi\in I}$, write
\[
 \Nphi(I)=
 \sup_{\|a\|_2\le1}
 \left\|\sum_{\xi\in I}a_\xi\varphi_\xi\right\|_\Phi.
\]
The source proves that every bounded orthogonal system of size $n$ has, with
positive probability and hence in particular deterministically, a subset
$I$ with
\[
 |I|\ge \frac{n}{e(\log n)^{\alpha+1}}
 \quad\text{and}\quad
 \Nphi(I)\le C_\alpha(\log\log n)^{\alpha/2}.
\]
It asks whether the factor $(\log\log n)^{\alpha/2}$ can be improved.

\section{Proof intuition}

The source lower bound finds a consecutive Fourier block inside a random
subset.  That mechanism cannot control a deterministic subset, which may
avoid all long intervals.  The replacement is to use all binary characters
and search for an affine subspace of frequencies.  An affine $r$-flat is an
ideal bad configuration: the sum of its $m=2^r$ characters is supported on
the annihilator of its direction space, an event of probability exactly
$1/m$, and has normalized height exactly $\sqrt m$ there.

It remains to force such a flat inside \emph{every} target-size subset.  If
$A\subset\F_2^d$ has density $\delta$, the normalized number of ordered
$r$-dimensional cubes in $A$ is at least $\delta^{2^r}$.  This follows by
iterating Cauchy--Schwarz (in the particularly simple form of Jensen's
inequality for opposite cube faces).  Degenerate cubes, whose directions are
linearly dependent, account for fewer than $2^r/2^d$ of all direction
tuples.  Hence $A$ contains a genuine affine $r$-flat as soon as
$\delta^{2^r}2^d>2^r$.

At the source density
$\delta\asymp(\log N)^{-(\alpha+1)}$, this condition permits
$2^r\asymp_\alpha\log N/\log\log N$.  Since the Orlicz norm of the resulting
spike grows as $(\log 2^r)^{\alpha/2}$, the desired endpoint loss follows.

\section{The deterministic lower theorem}

Let $G_d=\F_2^d$, let $N=|G_d|=2^d$, and give $G_d$ uniform probability
measure.  For $\xi,x\in G_d$, define
\[
 \chi_\xi(x)=(-1)^{\xi\cdot x},
\]
where the dot product is over $\F_2$.  The $N$ functions
$(\chi_\xi)_{\xi\in G_d}$ are orthonormal and have supremum norm one.

\begin{theorem}[Worst-case optimality]\label{thm:main}
Fix $\alpha>0$ and let $\Phi$ be the Young function above.  There are
$c_\alpha>0$ and $d_0(\alpha)$ such that for every $d\ge d_0(\alpha)$,
the character system $(\chi_\xi)_{\xi\in\F_2^d}$ has the following
property.  Every $I\subseteq\F_2^d$ satisfying
\[
 |I|\ge \frac{N}{e(\log N)^{\alpha+1}},
 \qquad N=2^d,
\]
obeys
\[
 \Nphi(I)\ge c_\alpha(\log\log N)^{\alpha/2}.
\]
Consequently the factor $(\log\log n)^{\alpha/2}$ in the source upper
bound is optimal, up to constants depending only on $\alpha$, in the
worst-case best-subsystem problem.
\end{theorem}

We first prove the combinatorial statement which replaces random block
selection.

\begin{lemma}[Dense sets contain affine cubes]\label{lem:cube}
Let $A\subseteq\F_2^d$ have density $\delta=|A|/2^d$.  If an integer
$r\ge1$ satisfies
\[
 \delta^{2^r}2^d>2^r,
\]
then $A$ contains an affine subspace of dimension $r$.
\end{lemma}

\begin{proof}
Put $G=\F_2^d$, $N=|G|$, and $f=\one_A$.  For $s\ge0$, let
\[
 C_s=
 \E_{x,h_1,\ldots,h_s\in G}
 \prod_{\omega\in\{0,1\}^s}
 f\left(x+\sum_{j=1}^s\omega_jh_j\right).
\]
Thus $C_s$ is the normalized number of ordered $s$-cubes contained in
$A$, with degenerate cubes included, and $C_0=\delta$.

For fixed $h_1,\ldots,h_s$, set
\[
 g(x)=\prod_{\omega\in\{0,1\}^s}
 f\left(x+\sum_{j=1}^s\omega_jh_j\right).
\]
Splitting an $(s+1)$-cube into its two opposite $s$-faces and averaging
first in its final direction gives
\[
 C_{s+1}
 =\E_{h_1,\ldots,h_s}
   \E_{x,t}g(x)g(x+t)
 =\E_{h_1,\ldots,h_s}(\E_xg(x))^2
 \ge C_s^2.
\]
The last step is Jensen's inequality.  Induction yields
\[
 C_r\ge\delta^{2^r}.
\]
Therefore the number of ordered tuples
$(x,h_1,\ldots,h_r)$ defining a cube in $A$ is at least
$\delta^{2^r}N^{r+1}$.

We now bound degenerate tuples.  Choosing the directions in order, if
$h_1,\ldots,h_r$ are linearly dependent, then for some first index $j$ one
has
$h_j\in\operatorname{span}(h_1,\ldots,h_{j-1})$.  For each $j$, there are
at most
\[
 N^{j-1}2^{j-1}N^{r-j}=2^{j-1}N^{r-1}
\]
such ordered direction tuples.  Hence fewer than
$(2^r-1)N^{r-1}$ direction tuples are dependent, and after choosing $x$
fewer than $(2^r-1)N^r$ ordered cubes are degenerate.

The assumed inequality makes the total number of cubes strictly larger
than the number of degenerate tuples.  Thus one cube has linearly independent
directions.  Its $2^r$ vertices form an affine $r$-dimensional subspace
contained in $A$.
\end{proof}

\begin{lemma}[The target density forces a large flat]\label{lem:scale}
Put $\rho=\alpha+1$.  For all sufficiently large $N=2^d$, every
$A\subseteq\F_2^d$ with
\[
 |A|/N\ge \frac1{e(\log N)^\rho}
\]
contains an affine subspace of cardinality $m=2^r$ satisfying
\[
 \frac{\log N}{16\rho\log\log N}<m
 \le \frac{\log N}{8\rho\log\log N}.
\]
\end{lemma}

\begin{proof}
Let $m=2^r$ be the largest power of two not exceeding
\[
 M=\frac{\log N}{8\rho\log\log N}.
\]
For large $N$, $m>M/2$, giving the displayed bounds.  If
$\delta=|A|/N$, then
\begin{align*}
 \log\left(\frac{\delta^mN}{m}\right)
 &\ge \log N-m-\rho m\log\log N-\log m\\
 &\ge \log N-M(1+\rho\log\log N)-\log M.
\end{align*}
For all sufficiently large $N$, the last expression is positive (indeed it
is at least $\tfrac34\log N-\log M$).  Thus
$\delta^mN>m$, and Lemma~\ref{lem:cube} applies.
\end{proof}

\begin{lemma}[An affine character flat is an Orlicz spike]
\label{lem:spike}
There are $b_\alpha>0$ and $m_0(\alpha)$ such that the following holds.
If $V=\xi_0+H\subseteq\F_2^d$ is an affine subspace of cardinality
$m\ge m_0(\alpha)$, then
\[
 \left\|m^{-1/2}\sum_{\xi\in V}\chi_\xi\right\|_\Phi
 \ge b_\alpha(\log m)^{\alpha/2}.
\]
\end{lemma}

\begin{proof}
Let
\[
 H^\perp=\{x\in\F_2^d:h\cdot x=0\text{ for every }h\in H\}.
\]
Character orthogonality on the group $H$ gives
\[
 \left|m^{-1/2}\sum_{\xi\in V}\chi_\xi(x)\right|
 =m^{-1/2}\left|\sum_{h\in H}\chi_h(x)\right|
 =\begin{cases}
 \sqrt m,&x\in H^\perp,\\
 0,&x\notin H^\perp.
 \end{cases}
\]
Since $|H^\perp|=N/m$, the nonzero event has probability $1/m$.

Set $\lambda=b_\alpha(\log m)^{\alpha/2}$.  For large $m$,
$\sqrt m/\lambda\ge u_0$ and
\[
 \log(\sqrt m/\lambda)\ge\tfrac14\log m.
\]
Consequently
\[
 \E\Phi\left(
 \frac{m^{-1/2}|\sum_{\xi\in V}\chi_\xi|}{\lambda}
 \right)
 =\frac1{\lambda^2}
 \log^\alpha(\sqrt m/\lambda)
 \ge\frac{4^{-\alpha}}{b_\alpha^2}.
\]
Choose, for example, $b_\alpha=2^{-\alpha-1}$.  The last quantity is at
least $4>1$.  By the definition of the Luxemburg norm, the norm is at least
$\lambda$.
\end{proof}

\begin{proof}[Proof of Theorem~\ref{thm:main}]
Let $I$ have the stated cardinality.  By Lemma~\ref{lem:scale}, $I$ contains
an affine subspace $V$ of cardinality $m$ with
\[
 m\asymp_\alpha\frac{\log N}{\log\log N}.
\]
Put coefficients $a_\xi=m^{-1/2}$ on $V$ and zero on $I\setminus V$.
They have $\ell^2$ norm one, so Lemma~\ref{lem:spike} gives
\[
 \Nphi(I)\ge b_\alpha(\log m)^{\alpha/2}.
\]
For sufficiently large $N$,
\[
 \log m
 \ge\log\log N-\log\log\log N-\log(16\rho)
 \ge\tfrac12\log\log N.
\]
Absorbing the factor $2^{-\alpha/2}$ into the constant proves the lower
bound.  The source upper bound holds for every uniformly bounded orthogonal
system, so the two estimates establish optimality in worst-case asymptotic
order.
\end{proof}

\section{Verification and novelty notes}

The proof has no unproved dependency and no computational step.  Its two
quantitative thresholds can be audited directly:
\begin{itemize}
\item the cube lower bound counts all ordered cubes, while the explicit
      $(2^r-1)N^{r-1}$ estimate counts every dependent direction tuple;
\item at density $e^{-1}(\log N)^{-\rho}$, taking
      $2^r\le \log N/[8\rho\log\log N]$ leaves a positive exponential
      margin in the condition $\delta^{2^r}N>2^r$;
\item the affine character sum is an exact two-valued function, so the
      Luxemburg modular calculation has no tail estimate or asymptotic
      approximation other than choosing $N$ sufficiently large.
\end{itemize}

A bounded novelty check on August 11, 2026 searched the run registry and
solution/attempt indexes for arXiv:2505.00155, endpoint Orlicz optimality,
Walsh systems, affine subspaces, and cube-count terms.  External searches
covered the exact title and arXiv id, combinations of ``bounded orthogonal
systems'', ``Orlicz'', ``Walsh'', ``affine subspace'', and ``Gowers norm'',
and recent papers returned for affine-subspace statistics.  The search found
the source paper and unrelated/background work, but no later paper claiming
to answer Remark 2.11 and no matching deterministic lower theorem.  Novelty
therefore remains subject to expert literature review.

\section{Limitations and human-review recommendation}

The theorem is stated on the infinite dimension subsequence $N=2^d$, which
is sufficient to prove that the source factor cannot be uniformly replaced
by a smaller asymptotic order.  Constants and the threshold $d_0$ depend on
$\alpha$, exactly as in the source upper bound.  No assertion is made about
the sharp numerical constant.

This is recommended for expert review as a likely valid full solution of the
source optimality question.  The highest-value checks are: (i) the
$C_{s+1}\ge C_s^2$ cube-count recursion, (ii) the subtraction of degenerate
direction tuples, (iii) the scale choice in Lemma~\ref{lem:scale}, and
(iv) the precise match between the source's best-subsystem quantifiers and
the universal statement for the constructed character system.

\begin{thebibliography}{9}
\bibitem{Burstein}
W. Burstein,
\emph{$\Lambda_p$ Style Bounds in Orlicz Spaces Close to $L^2$},
arXiv:2505.00155v2, 2025.
\end{thebibliography}

\end{document}
